
\section{Equilibrium and Welfare}

\begin{frame}{Producer and Consumer Surplus}
    \showmytikz[1]{../shared/tikz/tikz-demandandsupplysimple.tex}
    \mycomment{
        On the market, demand and supply meet: there is one quantity of goods $x^*$ traded at price $p^*$.
        
        There may be a price which is so high that nobody will want to buy anymore. (In reality, there are almost always some few very rich fools willing to pay any which price.)
    
        There is also a saturation quantity at price zero, where everyone who wants to have the good will have it.
    }    
\end{frame}

\begin{frame}{Producer and Consumer Surplus}
    \showmytikz[2]{../shared/tikz/tikz-demandandsupplysimple.tex}
    \mycomment{
        Some buyers are happy: they would be willing to pay more than they have to. They enjoy a surplus - free benefit, the consumer surplus. This equals the sum of the difference between their reservation prices and the equilibrium price.

        Think about this as money you save that you would be willing to spend on the good if you had to. 
        
        There are also some who would only buy at lower price. They are better off not trading at all. They remain outside the market.
    }    
\end{frame}

\begin{frame}{Producer and Consumer Surplus}
    \showmytikz[3]{../shared/tikz/tikz-demandandsupplysimple.tex}
    \mycomment{
        The sellers are happy: they would even sell at a lower price. They also enjoy a surplus, a free benefit, the \textbf{producer surplus.} If they face no fixed costs, then these are their profits.

        There are also some who would sell at a higher price. They are better off not trading at all. They remain outside the market.
 
    }    
\end{frame}

\begin{frame}{Producer and Consumer Surplus}
    \showmytikz[2,3,4]{../shared/tikz/tikz-demandandsupplysimple.tex}
    \mycomment{
        The sum of consumer and producer surplus is the \textbf{welfare gain of trade.}

        This is the case with \textbf{perfect competition.} the equilibrium reached is Pareto-optimal.
    }    
\end{frame}

\begin{frame}{Producer and Consumer Surplus}
    \mycomment{
    Consider a situation with highly  inelastic demand and elastic supply.

    The biggest surplus is the consumer surplus to not surprise. Consumers are willing to pay much more, but the do not have to.    

    This situation is petty typical when consumers have no or few substitutes. 
    }
    \showmytikz[1,2]{../shared/tikz/tikz-demandandsupplysimple-inelasticdemand}
    
\end{frame}

\begin{frame}{Producer and Consumer Surplus}
    \mycomment{
    If instead demand is elastic and supply inelastic, the main surplus goes to the producers. 
    
    inelastic supply is not very typical, but appears when there are high entrance barriers or natural monopolies. Even a very large increase in the price for paintings by Leonardo da Vinci will not bring many more such paintings on the market. But if the price for those paintings would fall steeply, many more would be willing to buy. Very few people would be willing to pay more for those paintings than today's price, so the consumer surplus is low.
    }

    \showmytikz[1,2]{../shared/tikz/tikz-demandandsupplysimple-inelasticsupply}
\end{frame}

